\chapter{Doppler Radar Systems}
\label{ch:radar}
% Retain historical bibliography records formerly cited by this chapter.
% Their inclusion does not support the corrected claims or validate current products.
\nocite{us12186643,flightscopex3,homeperformancelab,manzellaforum}

Radar observations constrain a moving target through frequency, phase, amplitude and timing.
A launch monitor combines those observations with geometry, signal processing and, for some outputs, a collision or flight model.
The useful question is which quantities the available observations identify, under which assumptions, and with what uncertainty.
Those distinctions matter when launch-monitor outputs are used to explain a golfer's swing: two reported parameters can share the same measured input and therefore share its errors.

\begin{keypoint}
\textbf{What This Shows:} The measurement equations below explain radial velocity, angular ambiguity, spin signatures and conditional trajectory-based spin-axis inference.
They also show how an inferred face angle can inherit uncertainty from ball launch and club path.
\textbf{What It Does Not Show:} A patent disclosure, vendor specification or manufactured calculation does not establish the accuracy or internal implementation of every commercial product.
Device generation, firmware, ball type, geometry and operating mode must accompany any performance claim.
\end{keypoint}

\section{Continuous-Wave Doppler Fundamentals}
\label{sec:cwdoppler}

A continuous-wave (CW) radar transmits a carrier at frequency $f_c$, with wavelength $\lambda=c/f_c$.
For monostatic transmission and reception, target speeds small compared with $c$, and a sign convention assigning positive Doppler to increasing range, the instantaneous range rate $v_r=\dot r$ gives
\begin{equation}
\label{eq:doppler}
f_d = \frac{2v_r}{\lambda}=\frac{2f_cv_r}{c}.
\end{equation}
The factor of two accounts for the outgoing-plus-return path; the displayed Doppler sign depends on the receiver's mixing convention.
At an illustrative carrier of \SI{24.125}{\giga\hertz}, one mph of radial speed gives approximately \SI{71.95}{\hertz}; at \SI{10.5}{\giga\hertz}, it gives \SI{31.31}{\hertz}, using $c=\SI{299792458}{\meter\per\second}$ and $1\,\mathrm{mph}=\SI{0.44704}{\meter\per\second}$.
These conversions do not identify a particular product's carrier or detection range.

Carrier frequency alone does not set a flight-tracking limit.
For a coherent observation of duration $T$, Fourier-bin spacing of order $1/T$ corresponds to radial-velocity spacing of order $\lambda/(2T)$.
That is spectral spacing, not a universal velocity-error bound: windowing, acceleration, estimator design and signal-to-noise ratio matter.
For an ideal wideband ranging waveform, range resolution instead scales with transmitted bandwidth $B$, approximately $c/(2B)$; angular performance depends on aperture and geometry.
Detection range additionally depends on transmitted power, antenna gain, target scattering, propagation, clutter and detection thresholds.
A blanket thirty-yard limit for consumer \SI{24}{\giga\hertz} radar cannot follow from the carrier frequency; the hybrid disclosure discussed below itself proposes much longer ranges~\cite{fullswingpatent}.

\section{From Velocity Sensor to 3D Tracker: Phase Interferometry}
\label{sec:interferometry}

A single unmodulated CW channel supplies radial-velocity information without independently resolving absolute range and two angular coordinates.
Multiple calibrated receive antennas can add direction information.
Let $\uvec u$ denote the unit direction from the array toward the target.
For a locally planar wavefront in a narrowband, far-field approximation, receivers separated by baseline $\vect{d}$ have a relative delay $\tau=(\vect{d}\cdot\uvec{u})/c$, with direction and phase signs chosen consistently.
For baseline length $d$, direction cosine $u=\vect d\cdot\uvec u/d$ along it and a calibrated measured phase $\Delta\varphi$, the possible directions are
\begin{equation}
\label{eq:interferometry}
\Delta\varphi \equiv \frac{2\pi d}{\lambda}u \pmod{2\pi},
\qquad
u_k=\frac{\lambda(\Delta\varphi+2\pi k)}{2\pi d},
\quad k\in\mathbb Z,\quad |u_k|\leq1.
\end{equation}
The integer branch is part of the inverse problem.
For $d=\lambda$, direction cosines $0.2$ and $-0.8$ give the same complex phase.
Half-wavelength spacing avoids interior aliases over the full direction-cosine interval, but its opposite endpoints still coincide in phase.
Larger baselines may be usable over a restricted field of view or with additional, appropriately chosen baselines and joint ambiguity resolution; an antenna count alone does not guarantee uniqueness~\cite{us9958527}.

Two independent baselines constrain two direction cosines, with a physical hemisphere choice needed to reconstruct a direction.
Near endfire, converting a direction cosine to an angle is poorly conditioned: for an angle $\theta$ measured from broadside so that $u=\sin\theta$, small unwrapped phase errors give $\delta\theta\simeq\lambda\delta\varphi/(2\pi d\cos\theta)$.
This local approximation excludes branch errors and becomes unreliable near $\cos\theta=0$.
Near-field wavefront curvature requires a range-dependent propagation model rather than the plane-wave formula.
Combining a resolved direction with independently estimated range and range rate can yield a three-dimensional track, provided the channels observe the same target and share calibrated coordinates and timing.
A per-Doppler-bin direction is a direction of contributing scattering, not automatically the ball center or a labeled point on a clubface~\cite{us10850179}.

\begin{implication}
One speed module and two angle modules form an investigational sensing arrangement, not a validated launch monitor merely by supplying three outputs.
RFbeam's K-LD7 exposes raw I/Q data: its RADC packet contains Rx1 and Rx2 at one carrier setting and Rx1 at a second setting used for FSK ranging.
The March 2021 Revision B datasheet specifies a maximum configured speed of \SI{100}{\kilo\meter\per\hour}, a typical \SI{29}{\milli\second} frame at that setting, and warnings about out-of-range returns and serial throughput~\cite{kld7}.
Its person-detection example and configurable range do not establish golf-ball performance.
An implementer must establish speed and signal bandwidth, target association, synchronization, ambiguity handling and error propagation before treating these modules as a three-dimensional ball tracker.
A rigid mount and surveyed calibration help, but do not remove the need to model drift, multipath and processing errors.
\end{implication}

\section{CW + FMCW Hybrids}
\label{sec:fmcw}

Frequency-modulated continuous-wave (FMCW) sensing adds delay information through a known transmitted frequency history.
For a moving target, beat frequency contains range and Doppler contributions, so separating them depends on the waveform and estimator; a chirp does not make range an assumption-free observable.
Full Swing's US\,11,311,789 describes example \SI{24}{\giga\hertz} CW and linear-FMCW channels, multiplexing alternatives and several antenna layouts~\cite{fullswingpatent}.
The disclosure associates Doppler signatures with speed and spin and modulated signals with range, and discusses ranges exceeding \SI{250}{\meter}.
Its examples include a ball swung on a line at \SI{75}{\meter}; that demonstration is not an independent accuracy study of full golf shots.
The array drawings include both symmetric and nonuniform arrangements, so one example should not be presented as the definitive commercial KIT layout.
The disclosure establishes proposed mechanisms; which mechanisms are deployed, and their error distributions, require separate product-specific evidence.

\section{Ball Measurement}

\subsection{Speed and Launch Angles}

A processor must associate the post-impact Doppler return with the ball before interpreting it as ball radial velocity.
Three-dimensional launch velocity also requires direction information from a fitted track or another sensing channel.
If $\alpha$ is the angle between target velocity and the radar line of sight, $v_r=\lVert\vect v\rVert\cos\alpha$; it is not generally enough to substitute the launch angle relative to the ground for $\alpha$.
Dividing by $\cos\alpha$ becomes poorly conditioned near a tangential trajectory, and uncertainty in direction then strongly affects inferred speed.
The reported impact state must also distinguish a sample taken after impact from a value extrapolated back to the defined event.

Leach et al. compared TrackMan Pro IIIe and Foresight GC2+HMT with a benchmark optical system using 240 shots across a driver, 7-iron and utility wedge~\cite{leach2017}.
They found generally stronger agreement for ball parameters than clubhead parameters, with several results dependent on club.
That evidence supports a device- and study-specific distinction; it does not establish best agreement for the same parameters in every study or newer product.

\subsection{Spin Rate: Harmonic Signatures and Their Ambiguities}
\label{sec:radarspin}

Tuxen's US\,8,845,442 describes extracting spin information from Doppler spectral structure~\cite{us8845442}.
In a rigid-ball point-feature model, let $\vect v_C$ be center velocity, $\vect\rho(t)$ the feature's offset from the center, and $\vect\omega$ the angular velocity, all expressed in one frame.
Using a common line-of-sight unit vector $\uvec n$ across the small target gives
\begin{equation}
\label{eq:radarfeature}
f_{d,\mathrm{feature}}(t)=\frac{2}{\lambda}
\uvec n\cdot\left[\vect v_C+\vect\omega\times\vect\rho(t)\right].
\end{equation}
The rotational term vanishes for spin parallel to that line of sight.
For constant spin and approximately fixed viewing direction, a feature at perpendicular distance $\rho_\perp$ from the spin axis has a sinusoidal projected velocity with amplitude $\lVert\vect\omega\rVert\rho_\perp\sin\beta$, where $\beta$ is the spin-axis/viewing angle.
The simpler amplitude $\rho_\perp\lVert\vect\omega\rVert$ requires a view perpendicular to the spin axis; $\rho_\perp$ equals the ball radius only for an equatorial surface feature.
Actual returns also depend on coherent scattering, visibility and rotating reflectivity, so this point-feature model is not a complete electromagnetic ball model.

Let $S$ denote spin rate in revolutions per minute.
Periodic modulation at rotation frequency $f_{\mathrm{spin}}=\lVert\vect\omega\rVert/(2\pi)$ can produce spectral lines
\begin{equation}
\label{eq:sidebands}
f_n=f_d+n f_{\mathrm{spin}},\quad n\in\mathbb Z,
\qquad S\,[\mathrm{rpm}]=60 f_{\mathrm{spin}}\,[\mathrm{Hz}].
\end{equation}
Observed peaks need not include every integer order, symmetric peaks need not have equal detectable amplitude, and some rotation aspects may yield weak modulation.
For example, peaks offset by 200 and 400 Hz can represent the first two orders of 200 Hz or the second and fourth orders of 100 Hz.
A reflector pattern that repeats after half a revolution can suppress odd harmonics relative to the mechanical rotation frequency.
Consequently, multiplying the spacing of any two visible peaks by sixty is not a generally valid spin estimator.

The patent describes tracking candidate traces, qualifying harmonics and assigning their orders before estimating rotation frequency~\cite{us8845442}.
Comb or cepstral processing can organize candidates, but cannot by itself resolve information absent from the observation.
Window duration, frequency drift, noise, clutter, scattering asymmetry and ball markings affect which candidates are distinguishable.
A harmonic spacing yields a candidate spin magnitude; it does not alone give a three-dimensional spin axis, though additional phase, amplitude, geometric or trajectory observations may supply axis information.

\subsection{Spin Axis: Conditional Trajectory Inversion}
\label{sec:spinaxis}

Trajectory inversion combines observations with a force model.
Let $\vect v_a=\vect v-\vect w$ be air-relative velocity for wind $\vect w$; let $\vect A$ denote ball acceleration in the ground-fixed inertial approximation, $\vect g$ gravitational acceleration, and $\vect D$ and $\vect L$ drag and lift accelerations, not forces.
Assume aerodynamic acceleration separates into drag parallel to $\vect v_a$ and lift perpendicular to it, with no omitted transverse force.
For nonzero airspeed,
\begin{equation}
\label{eq:liftextract}
\vect D=\frac{(\vect A-\vect g)\cdot\vect v_a}
{\lVert\vect v_a\rVert^2}\vect v_a,
\qquad
\vect L=\vect A-\vect g-\vect D.
\end{equation}
Physical drag additionally requires the fitted parallel component to oppose airspeed; a projection alone does not enforce that sign.
Unknown wind, omitted forces and differentiation errors can contaminate the inferred lift.
Acceleration estimated from position is especially sensitive to noise and the chosen temporal fit.

For a fixed spin-axis direction over the fitted interval and a Magnus model with $\vect L$ perpendicular to that axis, stacking the constraints at $N$ sample times gives
\begin{equation}
\label{eq:radaraxis}
\mathbf M\uvec{\omega}=\vect0,
\qquad
\mathbf M=\begin{bmatrix}\vect L_1^{\mathsf T}\\\vdots\\\vect L_N^{\mathsf T}\end{bmatrix},
\qquad \lVert\uvec{\omega}\rVert=1.
\end{equation}
In exact consistent data, rank two leaves one axis line, with two possible signs; rank one leaves a plane of candidate vectors and therefore a continuum of unit axes.
No lift leaves no directional information through this constraint.
The zero vector is always a solution of the unconstrained homogeneous system, which is why unit normalization is essential.
Noisy rows generally do not have an exact common nullspace: a weighted unit-norm fit and its conditioning are needed, rather than declaring uniqueness from a numerical rank alone.
The separation between the smallest singular values matters to axis stability.

A signed Magnus relation, such as $\vect L=\kappa\vect\omega\times\vect v_a$ with a stated positive effective aerodynamic coefficient $\kappa$, can help select the orientation and estimate spin components.
At one instant, adding spin parallel to $\vect v_a$ leaves this cross product unchanged.
Changing air-relative velocity over time can add information if the same spin vector and force model remain applicable; instantaneous insensitivity does not prove non-identifiability over every trajectory.
The patent's trajectory route must be interpreted with those model and geometry assumptions~\cite{us8845442}.
A longer observed flight often supplies more useful curvature, but does not guarantee accurate inference.
A short indoor segment can be poorly conditioned, especially with low lift or uncertain wind; both outdoor and indoor fitted axes remain estimates with different information and uncertainty.

\subsection{A Dimensional and Algebraic Check on a Delay Disclosure}

FlightScope's US\,10,775,492 describes a different spin-axis approach using spatially separated receiver pairs~\cite{us10775492}.
In that disclosure, $S_V$ and $S_H$ are vertical and horizontal receiver separations, $D$ is target distance, and $\Phi$ is the spin-axis angle inferred from those pairs.
Its displayed relations give $T_V=(S_V/D)\cos\Phi$ and $T_H=(S_H/D)\sin\Phi$, while calling the $T$ quantities time delays, and subsequently give $\arctan[S_HT_H/(S_VT_V)]$ for the angle.
With $S_V,S_H,D$ all lengths, those right-hand sides are dimensionless and require an additional time scale to represent delays.
Moreover, even treating the printed $T$ quantities as normalized values, elimination gives
\begin{equation}
\label{eq:radardelaycheck}
\Phi=\operatorname{atan2}\!\left(\frac{T_H}{S_H},\frac{T_V}{S_V}\right),
\qquad
\tan\Phi=\frac{S_VT_H}{S_HT_V},
\end{equation}
where the tangent form applies only when its denominator is nonzero and loses quadrant information.
For manufactured values $S_V=0.1$, $S_H=0.2$, $D=3$ in consistent length units and $\Phi=30\degs$, the displayed patent ratio returns approximately $66.59\degs$; the algebraic elimination returns $30\degs$.
This identifies an inconsistency in the displayed relationships, not a validated replacement for the full measurement method or an interpretation of patent scope.
A working delay estimator still needs timing definitions, feature correspondence, rotational geometry and experimental validation.

\subsection{Flight Observation, Extrapolation, and Ball Modes}

A trajectory tracked to landing can support measured-flight outcomes; a trajectory truncated by a screen requires a flight model for the unobserved part.
A normalized trajectory, computed for selected weather or ball conditions, must also be distinguished from the actual observed shot.
These are output modes, not an absolute division between radar brands.
TrackMan's published TM4 minimum indoor setup of \SI{4.7}{\meter} is \emph{device-to-net} distance, not \SI{4.7}{\meter} of ball flight~\cite{trackmanspecs}.
FlightScope's X3C instructions separately specify sensor-to-tee and ball-flight distances and identify marking/RCT requirements for indoor spin capture~\cite{flightscopex3csetup}.
Instructions for that generation should not silently be applied to every earlier X3 or Mevo.

Garmin's general R10 guidance describes standard-ball conditions under which spin is modeled rather than captured, including ball flight below \SI{20}{\meter} or ball speed below 90 mph~\cite{garminr10support}.
It must be read alongside the RCT mode: Garmin describes indoor spin capture with compatible Titleist RCT balls and current software, with at least eight feet tee-to-net and about two ball rotations for its most reliable results~\cite{garminr10rct}.
The support text allows successful capture with fewer rotations but also describes fallback to estimated spin when capture is inadequate; the app distinguishes that estimate typographically.
Thus the standard-ball twenty-meter/ninety-mph guidance is not a universal requirement to measure spin in every supported mode.
Marked radar-reflective balls and optically printed balls are different sensing aids and are not interchangeable across devices.

\section{Club Measurement}
\label{sec:radarclub}

\subsection{What the Radar Tracks on the Head}

A rotating clubhead has a velocity field: for a point $P$ offset by $\vect\rho$ from reference point $C$, $\vect v_P=\vect v_C+\vect\omega\times\vect\rho$.
Here $\vect\omega$ is clubhead angular velocity, rather than the ball angular velocity in the preceding section.
The fastest or strongest radar return need not correspond to the reference point used for reported club speed.
Its bias is not universally positive; it depends on geometry and which part of the head dominates the return.
TrackMan describes a reconstructed clubhead silhouette and defines reported parameters using specific points and event times~\cite{trackmanoert,trackmanclubspeed}.
Club path and attack angle concern a declared velocity direction; face angle and dynamic loft concern surface orientation, potentially at the impact location rather than the geometric center.
For curved faces, those orientations differ.
Point, frame, timing and surface definitions must therefore be reconciled before attributing a disagreement between instruments to sensor error.

\subsection{Face Angle and Dynamic Loft}
\label{sec:faceangle}

Radial-velocity sensing alone does not uniquely determine every clubface orientation, but it is too strong to say radar can never provide orientation information.
Resolved scattering geometry, multiple views, temporal constraints and other sensors can add information; the available model and measurements determine observability (\cref{app:screw}).
One possible estimation route is inversion of a local collision approximation.
Let $\varphi_f$ be face angle, $\varphi_p$ club path angle and $w_f$ the effective face weighting in a local collision model.
For small horizontal angles about the same reference direction, suppose launch direction satisfies $\varphi_{\mathrm{launch}}\simeq w_f\varphi_f+(1-w_f)\varphi_p$ with known nonzero $w_f$.
Then
\begin{equation}
\label{eq:faceinversion}
\varphi_f\simeq\frac{\varphi_{\mathrm{launch}}-(1-w_f)\varphi_p}{w_f}.
\end{equation}
This is an inverse of the stated local model, not a description of every radar product.
The effective weighting depends on collision conditions and must not be transferred unchanged between clubs, lofts, ball types and impact locations (\cref{sec:facepathweighting}).
A vertical-angle approximation likewise requires its own justified weighting and operating domain; it is not automatically obtained by substituting names into a horizontal fit.

For small perturbations within this model,
\begin{equation}
\label{eq:radarfaceerror}
\delta\varphi_f\simeq
\frac{\delta\varphi_{\mathrm{launch}}}{w_f}
-\frac{1-w_f}{w_f}\delta\varphi_p
+\frac{\varphi_p-\varphi_{\mathrm{launch}}}{w_f^2}\delta w_f.
\end{equation}
The launch-error gain is $1/w_f$ and therefore depends on the chosen collision weighting.
Correlated measurement errors and uncertainty in the collision model also contribute.
Off-center impact and gear effect can invalidate a fit that omits impact location; a richer estimator may use additional observations to address them.
A common horizontal reference offset moves face and path together and cancels in their difference in this approximation, whereas independent misalignments need not cancel.

\begin{keypoint}
An inferred face angle can be statistically coupled to the measured launch direction because the estimator uses that direction as an input.
Agreement or correlation between those outputs is therefore not independent confirmation of a physical face-to-launch relationship.
For swing analysis, document the measurement and inference dependencies, propagate their covariance, and compare against an independent criterion before drawing conclusions about technique.
\end{keypoint}

\subsection{Known Limitations}

The Leach et al. comparison motivates caution about clubhead outputs for research, but does not rank all modern radar and optical systems~\cite{leach2017}.
A pickup rate describes availability, not the error distribution of successful readings.
TrackMan's stated club-data pickup rate in its 2020 OERT description should be read in that sense~\cite{trackmanoert}.
Validation should retain missed shots, state the tested clubs and speeds, and distinguish bias, dispersion and repeatability from agreement with an independent reference.

\section{Optically Enhanced Radar: OERT and Fusion Tracking}
\label{sec:oert}

TrackMan's OERT description presents synchronized dual radar and camera observations, including clubhead silhouette and impact-location information~\cite{trackmanoert}.
Its current specification table describes TM4 and the indoor iO separately; the latter combines \SI{24}{\giga\hertz} radar with cameras and describes optical three-dimensional spin measurement~\cite{trackmanspecs}.
Those descriptions support the stated product distinctions without exposing every proprietary processing step.
FlightScope describes Fusion Tracking as synchronized image processing combined with radar, and US\,10,338,209 discloses mechanisms for combining the modalities~\cite{flightscopex3csetup,us10338209}.
The presence of a patent does not establish independent accuracy or prove that every disclosed option is deployed.

Rapsodo explicitly requires the printed RPT pattern for MLM2PRO optical spin measurement; its support distinguishes RPT from RCT radar-reflective balls~\cite{rapsodomlm2prorpt}.
Full Swing markets machine-learning radar and video capabilities, but that description alone does not establish that its camera contributes to a fused spin/face estimator~\cite{fullswingkit}.
A camera used for swing video is not automatically a measurement camera.

\begin{implication}
Choose an additional sensor to reduce a specified ambiguity or uncertainty, then test whether the combined estimator does so.
Optical features can constrain orientation or impact location; radar can contribute radial velocity and flight tracking; ball markings can improve identifiable rotational structure.
Their value depends on visibility, timing, calibration and the joint model.
Combining modalities can improve inference, but neither sensor count nor a vendor's architectural choice proves a universal limit of radar or the accuracy of a new implementation.
\end{implication}
